Los tres parámetros de esta sección no quedan sin cerrar por alcance o
descuido: cada uno tiene un veredicto explícito y una razón concreta por
la que la medición directa no es la vía adecuada hoy -- un permiso externo
pendiente en un caso, una decisión ética deliberada de no interferir con
infraestructura compartida en otro, y una dependencia real de resultados
que todavía no existen (Fase 4) en el tercero. Ninguno queda como
``trabajo futuro'' genérico: cada uno indica la condición exacta que, de
cumplirse, permitiría cerrarlo con medición.

\subsection{Niveles de frecuencia previstos (F0--F4)}

Los manifiestos de producción declaran cinco niveles fijos equiespaciados
(100\,\%, 75\,\%, 50\,\%, 25\,\% y 0\,\% del rango real de frecuencia
detectado en el nodo) además del nivel de referencia (gobernador nativo).
Esta elección no puede verificarse empíricamente hoy: requiere el permiso
de escritura de frecuencia (P1 para CPU, P4 para GPU), que sigue sin
llegar a la fecha de este informe. Tampoco se encontró en la literatura ya
revisada una recomendación explícita sobre cuántos puntos son suficientes
para caracterizar la curva DVFS de una carga: Guerreiro et al.
\citep{Guerreiro2019} infieren el resto del espacio de frecuencias a partir
de una única frecuencia base, sin barrer puntos uniformes; Ali et al.
\citep{Ali2023} parten igualmente de una configuración base y extrapolan de
forma analítica el resto del espacio, en vez de medir en un conjunto fijo
de puntos como aquí. La elección de cinco puntos equiespaciados es
razonable como decisión de diseño (visibiliza la forma de la curva sin
disparar el tamaño de la matriz experimental), pero no está anclada a un
estudio específico. Queda como pendiente explícito: en cuanto lleguen los
permisos, correr primero un reconocimiento con más puntos (por ejemplo,
nueve, cada 12{,}5\,\%) sobre dos o tres kernels representativos para
verificar si cinco puntos ya capturan la forma real de la curva.

\subsection{Umbral de carga externa (E08) y rango de temperatura (E02)}

\textbf{Corrección importante respecto a una versión anterior de este
análisis, y hallazgo más serio que una simple falta de barrido}: no es
solo que E08 y E02 se excluyan deliberadamente de prueba de estrés por
razones éticas -- una revisión del código de \texttt{preflight.py} y
\texttt{campaign.py} encontró que \textbf{ninguno de los dos chequeos se
ejecuta hoy en ninguna campaña real}. \texttt{run\_campaign\_preflight}
(el preflight que sí se invoca en producción) no incluye ni E02 ni E08;
ambos viven en \texttt{run\_reduced\_preflight}, una función con pruebas
unitarias propias pero que \texttt{campaign.py} nunca llama. Y aunque se
llamara, \texttt{Manifest} no declara el campo
\texttt{package\_temperature\_c} que \texttt{check\_temperature} necesita
-- con ese campo ausente (\texttt{None}), el chequeo aprueba de forma no
bloqueante con el mensaje ``no hay sensor disponible'', nunca bloquea
nada. El mismo tipo de hallazgo aplica al preflight de GPU (G01--G03): el
manifiesto de referencia GPU declara \texttt{gpu.enabled: false}, así que
esos chequeos tampoco corren hoy, y ni siquiera se les pasa un inspector
NVML desde \texttt{cli.py}. Este hallazgo se documenta con el mismo
detalle en la Sección~\ref{sec:metodologia} y en el balance final
(Sección~\ref{sec:balance}); aquí se completa con la justificación
teórica de cada umbral, ahora explícitamente como un umbral \emph{sin
mecanismo de verificación activo}, no como un umbral verificado pero sin
barrido de estrés.

E08 (carga externa normalizada $\leq 1{,}0$) se apoya en la convención
estándar de "nodo no saturado" en HPC, pero merece una precisión: el
propio \texttt{/proc/loadavg} cuenta procesos ejecutables o en espera de
E/S a nivel de \emph{todo el sistema}, no solo los de los núcleos
delegados a la corrida -- dividir esa cifra global entre 6 núcleos
delegados no aísla, por sí solo, si la contaminación ocurre específicamente
sobre esos 6 núcleos o en otra parte del nodo (irrelevante para la corrida
si el nodo tiene núcleos libres en otro socket). El umbral sigue siendo una
señal razonable de saturación general del nodo, pero no un chequeo preciso
de aislamiento por núcleo -- y, de cualquier forma, hoy no se ejecuta en
ninguna corrida real por el hallazgo anterior.

E02 (temperatura de paquete en $[0,90]\,^{\circ}\text{C}$) queda
\textbf{solo parcialmente} justificado por ficha técnica, no confirmado:
la hoja de especificaciones de Intel ARK para el SKU exacto del nodo (Xeon
Gold 5315Y, ARK 215286) documenta una temperatura de caja máxima
($T_{\text{case}}$) de 81\,$^{\circ}$C, no
90\,$^{\circ}$C.\footnote{\url{https://www.intel.com/content/www/us/en/products/sku/215286/}}
La temperatura de unión máxima ($T_j$máx, un límite distinto y típicamente
más alto que $T_{\text{case}}$) no tiene, hasta donde se pudo verificar en
esta revisión, una cifra pública específica para este SKU en ARK -- una
versión anterior de este análisis citaba un rango de 100--110\,$^{\circ}$C
sin fuente oficial que lo respaldara para este procesador en particular, y
se retira aquí por no estar verificado, en vez de mantenerlo sin cita. Lo
que sí queda establecido con fuente verificable es que $T_{\text{case}}$
(81\,$^{\circ}$C) es menor que el umbral configurado (90\,$^{\circ}$C),
así que el margen de seguridad no puede compararse directamente contra
$T_{\text{case}}$ sin aclarar que son magnitudes distintas: ninguna
corresponde exactamente a lo que \texttt{coretemp} reportaría si el sensor
estuviera conectado -- \texttt{coretemp} entrega un DTS (\textit{Digital
Thermal Sensor}) relativo a $T_j$máx, no una lectura directa de
$T_{\text{case}}$; y RAPL, aunque está habilitado en este nodo, mide
energía consumida, no temperatura -- ninguno de los dos provee hoy el dato
que \texttt{package\_temperature\_c} necesitaría. Diseñar E02 comparando
directamente una futura lectura DTS de \texttt{coretemp} contra
$T_{\text{case}}$ sería, por esta razón, metodológicamente incorrecto; el
chequeo, si se activa, debería compararse contra un límite expresado en
los mismos términos que el sensor real (DTS relativo a $T_j$máx). Así que
90\,$^{\circ}$C no puede describirse como "el umbral típico verificado", y
merece una precisión más incómoda que la de una versión anterior de este
análisis: 90\,$^{\circ}$C está \emph{por encima}, no por debajo, de la
única cifra con fuente oficial verificada para este SKU
($T_{\text{case}}=81\,^{\circ}\text{C}$). Que 90\,$^{\circ}$C sea o no un
margen conservador depende enteramente de $T_j$máx (típicamente mayor que
$T_{\text{case}}$ en procesadores Intel modernos, pero sin cifra pública
verificada para este SKU específico) -- una afirmación que no se puede
confirmar ni descartar con la evidencia disponible hoy. En ausencia de esa
cifra, lo único que puede decirse con evidencia real es que 90\,$^{\circ}$C
no está anclado a ningún límite térmico verificado y publicado por Intel
para este procesador exacto -- una versión previa de este análisis
afirmaba lo contrario (que 90\,$^{\circ}$C sí correspondía a un umbral
documentado), lo cual se corrige aquí.

\subsection{Peso del Producto Energía-Retardo (Fase 4)}

El peso $w$ de $EDP = E \cdot T^{w}$ tiene precedente para ambos valores
usuales en la literatura ya citada en el marco conceptual del proyecto --
$w=1$ (EDP clásico, Laros et al. \citep{Laros2013}) y $w=2$ ($ED^2P$,
evaluado junto con EDP por Ali et al. \citep{Ali2023} cuando se quiere
penalizar más fuertemente la degradación de rendimiento). Una versión
previa de este análisis proponía diferir la elección de $w$ hasta observar
los resultados de Fase 4 -- eso es metodológicamente incorrecto incluso sin
mala intención: decidir la métrica de evaluación \emph{después} de ver qué
resultado produce cada una abre la puerta, aunque no sea deliberado, a
escoger implícitamente la que favorezca al agente de control.

Este informe corrige eso fijando, desde ahora, $w=1$ (EDP clásico) como
\textbf{métrica primaria pre-registrada} de la Fase 4: es la formulación
con la que el objetivo del proyecto está planteado en el documento
metodológico principal, y no depende de ninguna medición pendiente. $ED^2P$
($w=2$) se reporta como \textbf{análisis de sensibilidad secundario} en
Fase 4 -- útil para mostrar si la conclusión cualitativa (el agente ahorra
energía sin degradar tiempo de forma inaceptable) es robusta a cuánto se
penaliza el retardo, pero nunca como la métrica que decide si el agente
"funciona". Fijar esto ahora, antes de tener datos de Fase 4, es
precisamente lo que evita la selección posterior de métrica.
